% Jalankan dari folder file ini dengan pdflatex + bibtex.
\documentclass[11pt,a4paper]{article}
\usepackage[utf8]{inputenc}
\usepackage[T1]{fontenc}
\usepackage{lmodern}
\usepackage[a4paper,margin=2.5cm]{geometry}
\usepackage{natbib}
\usepackage[hidelinks]{hyperref}
\title{Judul Artikel Ilmiah}
\author{Nama Penulis}
\date{\today}

\begin{document}
\maketitle
\begin{abstract}
[Tulis ringkasan sesuai evidence, bukan hasil yang belum tersedia.]
\end{abstract}

\section{Introduction}
[Latar masalah, gap penelitian, tujuan, dan kontribusi.]

\section{Related Work}
[Temuan sumber relevan dan perbedaan posisi studi.]

\section{Methodology}
[Metode yang telah dilakukan dan disetujui.]

\section{Results}
[Masukkan hanya hasil yang telah diverifikasi.]

\section{Discussion}
[Interpretasi, keterbatasan, dan implikasi.]

\section{Conclusion}
[Kesimpulan proporsional berdasarkan hasil.]

% Contoh referensi valid; hapus bila tidak relevan dengan artikel.
\nocite{lamport1994latex}
\bibliographystyle{plainnat}
\bibliography{../../../references/a7-pustaka}
\end{document}
